
\section{随机变量}
当我们把样本空间映射成数字来处理时，这时我们才有了更丰富的工具。

一个\textbf{随机变量}就是$\Omega\to \mathbb{R}$的一个函数，而且是可测的。这样就可以在实数轴上很好地研究它的性质了。

对于一个随机变量$X$，给出一个实数上的可测集合$S$，定义概率$P(X\in S)=P(\{e\in \Omega: X(e)\in S\})$。我们同样可以通过分离公理，用一个和$X$取值相关的谓词$\mathcal{P}(X)$描述集合$S=\{x\in \mathbb{R}:\mathcal{P}(x)\}$，这时我们简写$P(X\in S)$为$P(\mathcal{P}(X))$。
例如$P(1<X\leq 2)$表示$X$取值大于$1$小于等于$2$的概率。

随机变量有离散型和连续性之分，但本质上都可以用同一套理论描述：离散随机变量使用离散测度处理，连续随机变量使用Lebesgue测度处理。它们的结论基本上通过互换$\int$和$\Sigma$，$D$和$\Delta$就能互相转换。不过我在这里仍然使用不同的记号。

\subsection{随机变量的分布}

随机变量$X$的\textbf{累计分布函数}是这样一个函数$F_X(x)=P(X\leq x)$。很明显，它有这样的性质：
\[
\lim_{x\to -\infty} F_X(x)= 0, \lim_{x\to +\infty} F_X(x)= 1,\lim_{x\to x_0^+} F_X(x)= F(x_0)
\]

它们都是从前面概率的连续性，归一性推导出来的。

由概率的非负性，
\[
\forall x_1<x_2, F(x_2)-F(x_1)=P(x_1<X\leq x_2) \geq 0
\]
并且
\[
0\leq F(x)\leq 1
\]

从测度的视角看，概率密度函数 $f_X$ 是随机变量 $X$ 的分布测度 $P_X$ 相对于 Lebesgue 测度 $\lambda$ 的 Radon-Nikodym 导数。即：

若 $P_X$ 关于 Lebesgue 测度绝对连续（记为 $P_X \ll \lambda$），则存在非负可测函数 $f_X: \mathbb{R} \to [0, \infty)$，使得对任意可测集 $B \subseteq \mathbb{R}$，
\[
P_X(B) = \int_B f_X(t) \, d\lambda(t) = \int_B f_X(t) \, \dif t.
\]

此时 $f_X$ 称为 $X$ 的概率密度函数，它在几乎处处意义下唯一确定。下面给出的是更简单情况下的概率密度函数。

\subsubsection{离散型随机变量}

离散型随机变量只能取至多可数个值。我们可以通过列举所有可能取值发生的概率，来确定这个随机变量的统计规律。

如果一个离散型随机变量$X$可能能取值为$\{x_k\},k=1,2,\cdots$，则我们可以知道$P(X=x_k)$。我们知道，一个函数的纤维划分定义域，$\bigsqcup_{k=1}^{\infty}X^{-1}(x_k) = \Omega$。所以根据可数可加性：

\[
\sum_{k=1}^{\infty}P(X=x_k) =1
\]

虽然很简单，但是还是值得注意：$0\leq P(X=x_k)\leq 1$。

$P(X=x_k)=p_k$这种表现形式叫做随机变量的\textbf{分布律}，也叫\textbf{概率质量函数}。

\subsubsection{连续型随机变量}

描述连续型随机变量时，有概率密度函数这种更好的工具。对于一个离散型随机变量，它的累积分布函数一定是具有各种跳变不连续点的，并且在两个跳变点之间都是常值的。连续型随机变量的累计分布函数如果很大程度上（“几乎处处”）是可导的，此时也就可以用\textbf{概率密度函数}来描述。

一个连续型随机变量在某一点处的的概率密度函数是累计分布函数的导数（当然它也可能不存在）。有了概率密度函数$f_X$，我们就可以从积分来得出累积分布函数：

\[
F_X(x) = \int_{-\infty}^{x}f_X(t)\dif t
\]

它被称为概率密度函数也就是因为，$f_X(t)$表示了在这一处一个无限小的区间内的“概率质量”的集中程度。
$f_X(t)\dif t$就可以认为是$P(t<X<t+\dif t)$。

很明显，如果概率密度函数存在，那么它满足：

非负：
\[
f_X(x)\geq 0
\]
归一：
\[
\int_{-\infty}^{\infty}f_X(t)\dif t = 1
\]

事实上对于任意一个可积的非负函数，我都可以乘上一个系数让它归一，从而满足概率密度的要求。利用这一事实往往可以节约计算。

并且此时，我可以通过它求随机变量落在任意可测集合内的概率：

\[
P(X\in S)=\int_{S}f_X(t)\dif t
\]

特别地，$F_X(x)=P(X\leq x)=\int_{-\infty}^{x}f_X(t)\dif t$。

如刚才的叙述，它们在测度视角下都是一样的，只不过是换了种测度。

\subsubsection{随机变量的函数}

给一个随机变量$X$后复合上一个可测函数$\phi:\mathbb{R}\to\mathbb{R}$，我仍然会得到一个随机变量$\phi\circ X$。我们会把这个新的随机变量写成$\phi(X)$。

这个新的随机变量显然也存在自己的分布分布规律。它的累积分布函数也就可以表示成$F_{\phi(X)}(x)=P(\phi(X)\leq x)$。对于一般的连续型随机变量，想要直接通过概率密度函数来计算$\phi(X)$的概率密度函数是不太现实的，我们须要积分得到累积分布函数再找到$\phi(X)$的累计分布函数，再做微分。

不过当$\phi$是个严格单调函数时，就能直接计算出概率密度了。此时

\[
f_{\phi(X)}(y) = \frac{f_X(\phi^{-1}(y))}{|\phi'(\phi^{-1}(y))|} = f_X(\phi^{-1}(y))|(\phi^{-1})'(y)|
\]

我们在此不作证明，留在后面用一般性的结论处理。

作为一个简单的例子，数乘$X\mapsto\lambda X$是一个非常简单的随机变量函数。如果$\lambda=0$，情况非常简单，我们跳过。如果$\lambda \neq 0$，那么

\[
f_{\lambda X} (x) = \frac{1}{|\lambda|}f_X(\frac{x}{\lambda})
\]

对于离散型随机变量，$\phi(X)$的分布律可以通过对纤维直接求和得到：
\[
P(\phi(X)=y) = \sum_{x\in \phi^{-1}(y)}P(X=x)
\]
